1. ( \( \mathbf{1} \mathbf{~ p t}) \) Determine si el siguiente sistema de ecuaciones en \( \mathbb{R} \) tiene solución, si es así obtenga el espacio solución y si no, justifique la razón.
\[
\begin{aligned}
x_{1}-10 x_{2} & =-9 \\
x_{2}-10 x_{3} & =-9 \\
x_{3}-10 x_{4} & =-9 \\
x_{4}-10 x_{5} & =-9 \\
x_{5}-10 x_{6} & =-9 \\
x_{6} & =1
\end{aligned}
\]
\begin{solution}
Podemos escribir este sistema en forma matricial como:
\[
A \mathbf{x} = \mathbf{b}
\]
donde...
\[
A = \begin{bmatrix}
1 & -10 & 0 & 0 & 0 & 0 \\
0 & 1 & -10 & 0 & 0 & 0 \\
0 & 0 & 1 & -10 & 0 & 0 \\
0 & 0 & 0 & 1 & -10 & 0 \\
0 & 0 & 0 & 0 & 1 & -10 \\
0 & 0 & 0 & 0 & 0 & 1
\end{bmatrix}, \quad
\mathbf{x} = \begin{bmatrix}
x_{1} \\
x_{2} \\
x_{3} \\
x_{4} \\
x_{5} \\
x_{6}
\end{bmatrix}, \quad
\mathbf{b} = \begin{bmatrix}
-9 \\
-9 \\
-9 \\
-9 \\
-9 \\
1
\end{bmatrix}
\]
Notemos que todas las columnas son linealmente independientes por lo que la nulidad de A es cero, entonces la dimensión del dominio (6) es igual a su rango por Teorema de la Dimensión, esto nos es útil después para conocer cuántas soluciones hay. Para resolver \( A \mathbf{x} = \mathbf{b} \), aplicamos el método de eliminación gaussiana. Transformamos la matriz \( A \) en una matriz triangular superior mediante operaciones fila elementales. \\
\textit{Matriz aumentada inicial:}

\[
[A | \mathbf{b}] = \begin{bmatrix}
1 & -10 & 0 & 0 & 0 & 0 & \mid & -9 \\
0 & 1 & -10 & 0 & 0 & 0 & \mid & -9 \\
0 & 0 & 1 & -10 & 0 & 0 & \mid & -9 \\
0 & 0 & 0 & 1 & -10 & 0 & \mid & -9 \\
0 & 0 & 0 & 0 & 1 & -10 & \mid & -9 \\
0 & 0 & 0 & 0 & 0 & 1 & \mid & 1
\end{bmatrix}
\]

La matriz \( A \) es estrictamente triangular superior. Esta estructura indica que el sistema puede resolverse directamente hacia atrás sin necesidad de pivoteo así que realizamos esta sustitución. De la última ecuación, $x_6=1$. \\

Sustituimos hacia atrás:

\begin{enumerate}
    \item Ecuación 5: $x_5-10 \cdot 1=−9 \Rightarrow x_5=1$.
    \item Ecuación 4: $x_4-10 \cdot 1=−9 \Rightarrow x_4=1$.
    \item Ecuación 3: $x_3-10 \cdot 1=−9 \Rightarrow x_3=1$.
    \item Ecuación 2: $x_2-10 \cdot 1=−9 \Rightarrow x_2=1$.
    \item Ecuación 1: $x_1-10 \cdot 1=−9 \Rightarrow x_1=1$.
\end{enumerate}

Obtenemos la siguiente matriz aumentada resuelta.
\[
\begin{bmatrix}
1 & 0 & 0 & 0 & 0 & 0 & \mid & 1 \\
0 & 1 & 0 & 0 & 0 & 0 & \mid & 1 \\
0 & 0 & 1 & 0 & 0 & 0 & \mid & 1 \\
0 & 0 & 0 & 1 & 0 & 0 & \mid & 1 \\
0 & 0 & 0 & 0 & 1 & 0 & \mid & 1 \\
0 & 0 & 0 & 0 & 0 & 1 & \mid & 1
\end{bmatrix}
\]
De donde es directo ver que al estar compuesta por entradas canónicas su dominio e imagen son iguales, por Teorema de Rouché-Capelli se sabe que un sistema de ecuaciones lineales es consistente si y solo si el rango de la matriz de coeficientes es igual al rango de la matriz aumentada. Aquí ambas matrices tienen rango 6, asegurando la consistencia del sistema.\\
Para un sistema no homogéneo \(A \mathbf{x} = \mathbf{b}\), el espacio solución es un conjunto afín, que es una traslación del espacio solución del sistema homogéneo \(A \mathbf{x} = \mathbf{0}\) por la solución particular \(\mathbf{x}_p\).


\[ \therefore \textit{El espacio solución es  }
\mathbf{x} = \begin{bmatrix}
1 \\
1 \\
1 \\
1 \\
1 \\
1
\end{bmatrix} = span \begin{bmatrix}
0 \\
0 \\
0 \\
0 \\
0 \\
0
\end{bmatrix} + \begin{bmatrix}
1 \\
1 \\
1 \\
1 \\
1 \\
1
\end{bmatrix} 
\]
\end{solution}